\documentclass{article}[14pt]
\usepackage[T2A]{fontenc}
\usepackage{geometry}
\usepackage{fancyhdr}
\usepackage{tocloft}
\usepackage{hyperref}
\usepackage{titlesec}
\usepackage{graphicx}
\usepackage{caption}
\usepackage{wrapfig}
\usepackage[english, russian]{babel}
\geometry{
	a4paper,
	top=20mm, 
	right=25mm, 
	bottom=25mm, 
	left=25mm
}
\usepackage{amssymb}
\usepackage{fancyvrb}
\usepackage{fvextra}
\usepackage{ dsfont }
\usepackage{amsmath}
\usepackage{listings}
\usepackage{xcolor}
\usepackage{listingsutf8}
\usepackage{pdfpages}
\usepackage{lipsum}
\usepackage{makecell}
\usepackage{float}
\usepackage{longtable}
\usepackage{tabularx}


\lstset{ 
	basicstyle=\ttfamily\small,
	keywordstyle=\color{blue},
	inputencoding=utf8,              
	extendedchars=true,
	stringstyle=\color{red}
	commentstyle=\color{green},
	showspaces=false,         
	showstringspaces=false,  
	numbers=left,              
	numberstyle=\tiny\color{gray},
	breaklines=true,          
	frame=single,           
	captionpos=b,             
	tabsize=4,
	escapeinside={(*@}{@*)},
	literate={ё}{{\"o}}1 {Ё}{{\"O}}1 {«}{{\guillemotleft}}1 {»}{{\guillemotright}}1
	{А}{{\CYRA}}1 {Б}{{\CYRB}}1 {В}{{\CYRV}}1 {Г}{{\CYRG}}1 {Д}{{\CYRD}}1 
	{Е}{{\CYRE}}1 {Ж}{{\CYRZH}}1 {З}{{\CYRZ}}1 {И}{{\CYRI}}1 {Й}{{\CYRISHRT}}1 
	{К}{{\CYRK}}1 {Л}{{\CYRL}}1 {М}{{\CYRM}}1 {Н}{{\CYRN}}1 {О}{{\CYRO}}1 
	{П}{{\CYRP}}1 {Р}{{\CYRR}}1 {С}{{\CYRS}}1 {Т}{{\CYRT}}1 {У}{{\CYRU}}1 
	{Ф}{{\CYRF}}1 {Х}{{\CYRH}}1 {Ц}{{\CYRC}}1 {Ч}{{\CYRCH}}1 {Ш}{{\CYRSH}}1 
	{Щ}{{\CYRSHCH}}1 {Ъ}{{\CYRHRDSN}}1 {Ы}{{\CYRERY}}1 {Ь}{{\CYRSFTSN}}1 
	{Э}{{\CYREREV}}1 {Ю}{{\CYRYU}}1 {Я}{{\CYRYA}}1 
	{а}{{\cyra}}1 {б}{{\cyrb}}1 {в}{{\cyrv}}1 {г}{{\cyrg}}1 {д}{{\cyrd}}1 
	{е}{{\cyre}}1 {ж}{{\cyrzh}}1 {з}{{\cyrz}}1 {и}{{\cyri}}1 {й}{{\cyrishrt}}1 
	{к}{{\cyrk}}1 {л}{{\cyrl}}1 {м}{{\cyrm}}1 {н}{{\cyrn}}1 {о}{{\cyro}}1 
	{п}{{\cyrp}}1 {р}{{\cyrr}}1 {с}{{\cyrs}}1 {т}{{\cyrt}}1 {у}{{\cyru}}1 
	{ф}{{\cyrf}}1 {х}{{\cyrh}}1 {ц}{{\cyrc}}1 {ч}{{\cyrch}}1 {ш}{{\cyrsh}}1 
	{щ}{{\cyrshch}}1 {ъ}{{\cyrhrdsn}}1 {ы}{{\cyrery}}1 {ь}{{\cyrsftsn}}1 
	{э}{{\cyrerev}}1 {ю}{{\cyryu}}1 {я}{{\cyrya}}1               
}
\begin{document}
	
\thispagestyle{empty}
\begin{center}
	\Large {\bfseries МИНИСТЕРСТВО НАУКИ И ВЫСШЕГО ОБРАЗОВАНИЯ РОССИЙСКОЙ ФЕДЕРАЦИИ}
	\\
	\vspace{3mm}
	Федеральное государственное автономное образовательное учреждение высшего образования <<Санкт-Петербургский политехнический университет Петра Великого>>
	\\
	\vspace{3mm}
	Институт компьютерных наук и кибербезопасности
	\\
	\vspace{3mm}
	Высшая школа технологий искусственного интеллекта
	\\
	\vspace{3mm}
	Направление: 02.03.01 Математика и компьютерные науки
	\\
	\vspace{2cm}
	\large{
		Методы тестирования программного обеспечения}
	\\
	\LARGE{
		Отчет по лабораторной работе №2\\
		<<Тестирование методом черного и белого ящика>>}
\end{center}

\vspace{5.5cm}
\begin{flushleft}
	\Large
	Студент
	\\
	группы 5130201/30101
	\\
	\vspace{2cm}
	\Large
	Преподаватель
	\\ 
\end{flushleft}
\vspace{-4.6cm}
\begin{flushright}
	\Large
	\line(1,0){2cm}
	\hspace{0.4cm}
	Шишмаков В.О.
	\\
	\vspace{2,7cm}
	\hspace{2cm}
	\Large
	\line(1,0){2cm}
	\hspace{0.4cm}
	Курочкин М. А.
	\hspace{-0.1cm}
	
\end{flushright}

\vspace{1cm}
\hspace{9.8cm}
{
	\Large
	\vspace{0.8cm}
	<<\underline{\hspace{1cm}}>>\underline{\hspace{2.5cm}} 2026 г.
}
\hfill \break
\hfill \break
\hfill \break
\hfill \break
\begin{center} \large{Санкт-Петербург, 2026} \end{center}
\pagebreak

\titleformat{\section}
{\Huge\bfseries}
{\thesection\;}
{0.65cm}
{}
\titleformat{\subsection}
{\huge\bfseries}
{\thesubsection\;}
{0.65cm}
{}
\titleformat{\subsubsection}
{\Large\bfseries}
{\thesubsubsection\;}
{0.65cm}
{}
\titlespacing*{\section}
{0em}
{1.5ex plus 1ex minus 0.2ex}
{1.5ex plus 1ex minus 0.2ex}	

\renewcommand*\contentsname{\huge{Содержание}}
\Large{
	\tableofcontents}
\pagebreak

\section*{Введение}
\addcontentsline{toc}{section}{Введение}

Одним из важнейших этапов в процессе создания современного программного обеспечения является тестирование. Выделяют 
различные виды тестирования: модульное, интеграционное, функциональное, системное и приемочное. В рамках данной работы 
основное внимание уделяется модульному тестированию.

Модульное тестирование представляет собой базовый вид проверки, с которого, как правило, начинается 
оценка корректности функционирования программы. Оно предполагает проверку отдельных компонентов программного продукта: 
классов, модулей и функций. Основная цель модульного тестирования --- обнаружение дефектов, то есть случаев, когда реальное 
поведение модуля расходится с его спецификацией.

В практике модульного тестирования применяются два основных подхода: метод «черного ящика» и метод «белого ящика».

Оба этих подхода будут рассмотрены и применены к обеим программам в данной работе.

\pagebreak

\section{Постановка задачи}

\noindent \textbf{Цель:} провести тестирование двух программ методами черного и белого ящика.

\noindent \textbf{Задачи:}

\begin{enumerate}
    \item Ознакомиться с подходами к модульному тестированию, а именно с методами черного и белого ящика.
    \item Составить тесты для двух программ обоими методами и провести их.
    \item Провести анализ результатов тестирования.
\end{enumerate}

\pagebreak

\section{Описание методов тестирования}

\subsection{Метод черного ящика}

Одним из ключевых подходов к тестированию является метод «черного ящика», известный также как тестирование, 
ориентированное на данные, либо тестирование на основе входных и выходных параметров. Согласно данному подходу, 
программа рассматривается как <<черный ящик>>, внутренности и логика которого остаются скрыты и не 
принимаются во внимание. Основное внимание уделяется выявлению ситуаций, в которых реальная работа программы отклоняется 
от требований спецификации.

В рамках данного подхода необходимо выполнить проверку программы, используя такие техники, как разбиение входных данных на 
классы эквивалентности, анализ граничных значений и построение причинно-следственных диаграмм.

\subsubsection{Эквивалентное разбиение}

Качественно разработанный тест должен охватывать значительную часть прочих потенциальных тестов. В связи с этим 
целесообразно разделить все множество входных данных программы на ограниченное количество классов эквивалентности. 
Различают два типа классов: корректные классы, включающие допустимые входные значения, и некорректные классы, 
представляющие все остальные возможные состояния.

\subsubsection{Анализ граничных значений}

Граничные условия представляют собой ситуации, которые возникают в окрестностях предельных значений, определяющих 
классы эквивалентности. В отличие от эквивалентного разбиения, анализ граничных значений требует выбора элементов, 
позволяющих проверить именно каждую границу класса, и уделяет внимание как входным, так и выходным данным.

\subsubsection{Причинно-следственные диаграммы}

Причинно-следственная диаграмма --- формализованный графический способ описания спецификации. Каждый узел принимает 
значение 0 или 1. Данный метод восполняет недостаток подходов, основанных на анализе граничных значений и разбиении 
на классы эквивалентности, не учитывающих комбинации входных условий.

\subsection{Метод белого ящика}

В рамках подхода <<белого ящика>> предоставляется возможность изучать внутреннее устройство программы. 
Специалист по тестированию выбирает входные данные, анализируя логическую структуру программы.

\subsubsection{Покрытие операторов}

Минимальным требованием является выполнение каждого оператора программы как минимум однократно.

\subsubsection{Покрытие решений}

Необходимо подобрать такой набор тестов, при котором каждое логическое условие хотя бы один раз принимает 
значение <<истина>> и хотя бы один раз --- <<ложь>>.

\subsubsection{Покрытие условий}

Этот критерий требует, чтобы набор тестов обеспечивал для каждого элементарного условия, входящего в состав 
составного условия, выполнение как минимум один раз значения <<истина>> и один раз значения <<ложь>>.

\subsubsection{Покрытие решений и условий}

Набор тестов считается достаточным, если каждое элементарное условие принимает оба возможных значения, каждый 
возможный результат решения достигается минимум один раз и каждая точка входа активируется хотя бы однократно.

\subsubsection{Комбинаторное покрытие условий}

Данный критерий предполагает формирование такого набора тестов, который обеспечивает проверку всех возможных 
комбинаций значений элементарных условий в каждом логическом решении.

\pagebreak

%=============================================================
\section{Тестирование программы №1}
%=============================================================

\subsection{Описание программы}

\noindent \textbf{Автор:} Митяев Олег.

\noindent \textbf{Название:} Проверка простоты числа.

\noindent \textbf{Дано:}

\begin{itemize}
	\item $n$ --- целое число.
\end{itemize}

\noindent \textbf{Требуется:} определить, является ли число $n$ простым.

\noindent \textbf{Ограничения:}

\begin{enumerate}
	\item $n$ вводится как целое число.
	\item При некорректном вводе выводится сообщение \texttt{ошибка ввода} и выполняется повтор ввода.
	\item Числа $n \leq 1$ считаются непростыми.
	\item Проверка простоты выполняется перебором делителей до $[\sqrt{n}]$.
\end{enumerate}

\noindent \textbf{Спецификация:}

\begin{table}[h]
\centering
\begin{tabular}{|p{2.5cm}|p{3.5cm}|p{2.5cm}|p{4cm}|}
\hline
\textbf{Входные данные} & \textbf{Ожидаемый результат} & \textbf{Статус} & \textbf{Комментарий} \\
\hline
7 & вывод: простое & Корректно & Обычный случай (простое число). \\
\hline
4 & вывод: не простое & Корректно & Обычный случай (составное число). \\
\hline
2 & вывод: простое & Корректно & Граница: минимальное простое. \\
\hline
1 & вывод: не простое & Корректно & Граница: $n \leq 1$ непростое. \\
\hline
-5 & вывод: не простое & Корректно & Отрицательные числа непростые. \\
\hline
abc & вывод: ошибка ввода; ожидание повторного ввода & Ошибка ввода & Проверка обработки некорректного ввода. \\
\hline
80 12 & вывод: ошибка ввода; ожидание повторного ввода & Ошибка ввода & Обработка некорректного ввода (через пробел). \\
\hline
3.45 & вывод: ошибка ввода; ожидание повторного ввода & Ошибка ввода & Обработка некорректного ввода (не целое). \\
\hline
\end{tabular}
\end{table}

\begin{figure}[H]
	\centering
	\includegraphics[width=0.8\textwidth]{pics/scheme1.png}
	\captionof{figure}{Блок-схема программы №1}
\end{figure}

%-------------------------------------------------------------
\subsection{Тестирование методом <<черного>> ящика}
%-------------------------------------------------------------

\subsubsection{Разбиение на классы эквивалентности}

Составлено разбиение на классы эквивалентности, исходя из ограничений для тестируемой программы.

{\centering\large{
\begin{longtable}{|p{4.5cm}|p{5.5cm}|p{5.5cm}|}
	\caption{Классы эквивалентности для n}\\
	\hline
	Входное условие & Допустимые классы & Недопустимые классы \\ \hline
	Целое число $n$ & Отрицательное число (1) & Нецелое вещественное число (8) \\ \hline
	& Ноль (2) & Строка (буквы) (9) \\ \hline
	& Единица (3) & Несколько чисел через пробел (10) \\ \hline
	& Минимальное простое (4) & Пустая строка (11) \\ \hline
	& Четное число $>2$ (5) & Спецсимволы / смешанный ввод (12) \\ \hline
	& Нечетное простое число (6) & \\ \hline
	& Нечетное составное число (7) & \\ \hline
\end{longtable}}}

В таблице \ref{tab:test-black-box-equivalents} представлены тесты, покрывающие все классы эквивалентности.

{\centering\large{
\begin{longtable}{|p{0.4cm}|p{0.8cm}|p{3cm}|p{3.5cm}|p{3.5cm}|p{2cm}|}
	\caption{Тесты, покрывающие классы эквивалентности}
	\label{tab:test-black-box-equivalents}\\
	\hline
	N & n & Класс эквивалентности & Ожидаемый результат & Фактический результат & Результат \\ \hline
	1 & -5 & (1) & Число не простое & Число не простое & Пройден \\ \hline
	2 & 0 & (2) & Число не простое & Число не простое & Пройден \\ \hline
	3 & 1 & (3) & Число не простое & Число не простое & Пройден \\ \hline
	4 & 2 & (4) & Число простое & Число простое & Пройден \\ \hline
	5 & 8 & (5) & Число не простое & Число не простое & Пройден \\ \hline
	6 & 17 & (6) & Число простое & Число простое & Пройден \\ \hline
	7 & 15 & (7) & Число не простое & Число не простое & Пройден \\ \hline
	8 & 7.5 & (8) & Ошибка ввода; повторный ввод & Ошибка ввода; повторный ввод & Пройден \\ \hline
	9 & abc & (9) & Ошибка ввода; повторный ввод & Ошибка ввода; повторный ввод & Пройден \\ \hline
	10 & 12 5 & (10) & Ошибка ввода; повторный ввод & Ошибка ввода; повторный ввод & Пройден \\ \hline
	11 & <<>> & (11) & Ошибка ввода; повторный ввод & Ошибка ввода; повторный ввод & Пройден \\ \hline
	12 & 10a? & (12) & Ошибка ввода; повторный ввод & Ошибка ввода; повторный ввод & Пройден \\ \hline
\end{longtable}}}

Все 12 тестов прошли успешно.

\subsubsection{Анализ граничных условий}

Исходя из ограничений данной программы выделяем следующие граничные значения с точки зрения алгоритма:

\begin{enumerate}
	\item Переход от наибольшего непростого к наименьшему простому (1, 2).
	\item Переход между наименьшим четным составным (3, 4, 5).
	\item Делитель ровно на границе $\sqrt{n}$: $\sqrt{9} = 3$ (8, 9, 10).
\end{enumerate}

{\centering\large{
	\begin{longtable}{|p{0.4cm}|p{0.8cm}|p{4.5cm}|p{4.5cm}|p{2cm}|}
		\caption{Проверка граничных условий}\\
		\hline
		N & n & Ожидаемый результат & Фактический результат & Результат \\ \hline
		1 & 1 & Число не простое & Число не простое & Пройден \\ \hline
		2 & 2 & Число простое & Число простое & Пройден \\ \hline
		3 & 3 & Число простое & Число простое & Пройден \\ \hline
		4 & 4 & Число не простое & Число не простое & Пройден \\ \hline
		5 & 5 & Число простое & Число простое & Пройден \\ \hline
		6 & 8 & Число не простое & Число не простое & Пройден \\ \hline
		7 & 9 & Число не простое & Число не простое & Пройден \\ \hline
		8 & 10 & Число не простое & Число не простое & Пройден \\ \hline
	\end{longtable}
}}

При проведении тестирования программа успешно обработала все граничные условия.

\subsubsection{Причинно-следственная диаграмма}

Для построения причинно-следственной диаграммы введём следующие обозначения.

\textbf{Причины:}
\begin{enumerate}
    \item Ввод является целым числом.
    \item Число $n \leq 1$.
    \item Число $n = 2$.
    \item Число $n$ чётное и $n \neq 2$.
    \item Среди нечётных чисел от $3$ до $\lfloor\sqrt{n}\rfloor$ найден делитель $n$.
\end{enumerate}

\textbf{Следствия:}
\begin{enumerate}
    \setcounter{enumi}{20}
    \item Программа выводит сообщение об ошибке ввода и запрашивает повторный ввод.
    \item Программа выводит результат: \texttt{не простое} (по условию $n \leq 1$).
    \item Программа выводит результат: \texttt{простое} (граничный случай $n = 2$).
    \item Программа выводит результат: \texttt{не простое} (чётное число $n \neq 2$).
    \item Программа выводит результат: \texttt{не простое} (найден делитель при переборе).
    \item Программа выводит результат: \texttt{простое} (делитель не найден).
\end{enumerate}

\textbf{Промежуточные узлы:}
\begin{enumerate}
    \setcounter{enumi}{10}
    \item Причина 1 истинна И причина 2 ложна (корректный ввод, $n > 1$).
    \item Узел 11 истинен И причина 3 ложна И причина 4 ложна (число нечётное $> 2$, выполняется перебор).
\end{enumerate}

\begin{table}[H]
\centering
\caption{Таблица решений (метод чёрного ящика)}
\renewcommand{\arraystretch}{1.4}
\begin{tabular}{|l|c|c|c|c|c|c|c|}
\hline
 & \textbf{1} & \textbf{2} & \textbf{3} & \textbf{4} & \textbf{5} & \textbf{6} & \textbf{7} \\
\hline
Причина 1 (целое число)        & 0 & 1 & 1 & 1 & 1 & 1 & 1 \\
\hline
Причина 2 ($n \leq 1$)         & -- & 1 & 0 & 0 & 0 & 0 & 0 \\
\hline
Причина 3 ($n = 2$)            & -- & -- & -- & 1 & 0 & 0 & 0 \\
\hline
Причина 4 ($n$ чётное, $\neq2$) & -- & -- & -- & -- & 1 & 0 & 0 \\
\hline
Причина 5 (делитель найден)    & -- & -- & -- & -- & -- & 1 & 0 \\
\hline
Узел 11 ($C_1 \wedge \neg C_2$)             & 0 & 0 & 1 & 1 & 1 & 1 & 1 \\
\hline
Узел 12 ($U_{11} \wedge \neg C_3 \wedge \neg C_4$) & 0 & 0 & 0 & 0 & 0 & 1 & 1 \\
\hline
Следствие 21 (ошибка ввода)    & 1 & 0 & 0 & 0 & 0 & 0 & 0 \\
\hline
Следствие 22 (\texttt{не простое}, $n\leq1$) & 0 & 1 & 0 & 0 & 0 & 0 & 0 \\
\hline
Следствие 23 (\texttt{простое}, $n=2$)      & 0 & 0 & 0 & 1 & 0 & 0 & 0 \\
\hline
Следствие 24 (\texttt{не простое}, чётное)  & 0 & 0 & 0 & 0 & 1 & 0 & 0 \\
\hline
Следствие 25 (\texttt{не простое}, делитель найден) & 0 & 0 & 0 & 0 & 0 & 1 & 0 \\
\hline
Следствие 26 (\texttt{простое}, делитель не найден) & 0 & 0 & 0 & 0 & 0 & 0 & 1 \\
\hline
\end{tabular}
\end{table}

\begin{table}[H]
\centering
\caption{Тесты на основе таблицы решений}
\renewcommand{\arraystretch}{1.5}
\begin{tabular}{|c|p{4.2cm}|c|p{3.5cm}|p{3.5cm}|c|}
\hline
\textbf{N} & \textbf{Комбинация} & \textbf{n} & \textbf{Ожидаемый рез.} & \textbf{Фактич. рез.} & \textbf{Результат} \\
\hline
1 & $C_1=0$ & \texttt{abc} & Следствие 21: ошибка ввода & Следствие 21: ошибка ввода & Пройден \\
\hline
2 & $C_1=1,\ C_2=1$ & \texttt{-3} & Следствие 22: \texttt{не простое} & Следствие 22: \texttt{не простое} & Пройден \\
\hline
3 & $C_1=1,\ C_2=0,\ C_3=1$ & \texttt{2} & Следствие 23: \texttt{простое} & Следствие 23: \texttt{простое} & Пройден \\
\hline
4 & $C_1=1,\ C_2=0,\ C_4=1$ & \texttt{8} & Следствие 24: \texttt{не простое} & Следствие 24: \texttt{не простое} & Пройден \\
\hline
5 & $C_1=1,\ C_5=1$ & \texttt{15} & Следствие 25: \texttt{не простое} & Следствие 25: \texttt{не простое} & Пройден \\
\hline
6 & $C_1=1,\ C_5=0$ & \texttt{17} & Следствие 26: \texttt{простое} & Следствие 26: \texttt{простое} & Пройден \\
\hline
\end{tabular}
\end{table}

\subsubsection{Результаты тестирования программы №1 методом чёрного ящика}

В результате тестирования методом чёрного ящика было составлено 26 тестов
(12 классов эквивалентности + 8 граничных значений + 6 тестов по таблице решений). 
Все тесты были успешно пройдены, дефектов обнаружено не было.

%-------------------------------------------------------------
\subsection{Тестирование методом <<белого>> ящика}
%-------------------------------------------------------------

Алгоритм составления тестов методом «белого» ящика предполагает обход всех возможных путей
в теле программы и проверку выполнения каждого оператора не менее одного раза.

В программе можно выделить следующие ключевые условия:

\begin{enumerate}
    \item Условие C1: ввод является корректным целым числом.

        Ветвь F (False): вывод ошибки ввода, повторный запрос.

        Ветвь T (True): продолжение выполнения.

    \item Условие C2: $n \leq 1$.

        Ветвь T (True): вывод <<не простое>>, завершение.

        Ветвь F (False): продолжение.

    \item Условие C3: $n = 2$.

        Ветвь T (True): вывод <<простое>>, завершение.

        Ветвь F (False): продолжение.

    \item Условие C4: $n$ чётное и $n \neq 2$.

        Ветвь T (True): вывод <<не простое>>, завершение.

        Ветвь F (False): переход к перебору делителей.

    \item Условие C5: в теле цикла $i \leq \lfloor\sqrt{n}\rfloor$ (условие продолжения цикла).

        Ветвь T (True): продолжение итераций.

        Ветвь F (False): выход из цикла.

    \item Условие C6: $n \bmod i = 0$ (найден делитель).

        Ветвь T (True): вывод <<не простое>>, завершение.

        Ветвь F (False): переход к следующей итерации.
\end{enumerate}

\noindent \textbf{Возможные пути выполнения:}

\begin{enumerate}
    \item Некорректный ввод: C1(F) $\to$ повторный ввод.
    \item $n \leq 1$: C1(T) $\to$ C2(T) $\to$ <<не простое>>.
    \item $n = 2$: C1(T) $\to$ C2(F) $\to$ C3(T) $\to$ <<простое>>.
    \item Чётное $n > 2$: C1(T) $\to$ C2(F) $\to$ C3(F) $\to$ C4(T) $\to$ <<не простое>>.
    \item Нечётное составное: C1(T) $\to$ \ldots $\to$ C4(F) $\to$ цикл $\to$ C5(T) $\to$ C6(T) $\to$ <<не простое>>.
    \item Нечётное простое: C1(T) $\to$ \ldots $\to$ C4(F) $\to$ цикл $\to$ C5(T) $\to$ C6(F) повторяется $\to$ C5(F) $\to$ <<простое>>.
\end{enumerate}

\subsubsection{Покрытие операторов}

Для выполнения каждого оператора хотя бы один раз необходим тест, проходящий по основному пути
(нечётное простое число), при котором все операторы цикла выполняются.

\begin{table}[H]
\centering
\caption{Покрытие операторов}
\begin{tabularx}{\textwidth}{|c|c|X|X|c|}
\hline
\textbf{N} & \textbf{n} & \textbf{Ожидаемый результат} & \textbf{Фактический результат} & \textbf{Статус} \\
\hline
1 & 17 & простое & простое & Пройден \\
\hline
\end{tabularx}
\end{table}

Тест №1 обеспечивает выполнение всех операторов программы: ввод, проверки C2--C4, цикл с проверкой C5 и C6, вывод результата.

\subsubsection{Покрытие решений}

Необходимо, чтобы каждое из условий C1--C6 приняло оба значения (True и False).

\begin{table}[H]
\centering
\caption{Покрытие решений}
\begin{tabularx}{\textwidth}{|c|c|X|X|c|c|c|c|c|c|c|}
\hline
\textbf{N} & \textbf{n} & \textbf{Ожид. рез.} & \textbf{Факт. рез.} & \textbf{C1} & \textbf{C2} & \textbf{C3} & \textbf{C4} & \textbf{C5} & \textbf{C6} & \textbf{Статус} \\
\hline
1 & 17 & простое & простое & T & F & F & F & T/F & F & Пройден \\
\hline
2 & \texttt{abc} & ошибка ввода & ошибка ввода & F & -- & -- & -- & -- & -- & Пройден \\
\hline
3 & 0 & не простое & не простое & T & T & -- & -- & -- & -- & Пройден \\
\hline
4 & 2 & простое & простое & T & F & T & -- & -- & -- & Пройден \\
\hline
5 & 8 & не простое & не простое & T & F & F & T & -- & -- & Пройден \\
\hline
6 & 15 & не простое & не простое & T & F & F & F & T & T & Пройден \\
\hline
\end{tabularx}
\end{table}

Тесты №1 и №6 покрывают оба значения C5 и C6. Все условия протестированы на T и F.

\subsubsection{Покрытие условий}

Каждое элементарное условие должно принять значения True и False хотя бы по одному разу.
Условия C1--C6 являются элементарными. Набор тестов из предыдущего раздела полностью
удовлетворяет данному критерию.

\begin{table}[H]
\centering
\caption{Покрытие условий}
\begin{tabularx}{\textwidth}{|c|c|X|X|c|c|c|}
\hline
\textbf{N} & \textbf{n} & \textbf{Ожид. рез.} & \textbf{Факт. рез.} & \textbf{C1} & \textbf{C2} & \textbf{Статус} \\
\hline
1 & 17 & простое & простое & T & F & Пройден \\
\hline
2 & \texttt{abc} & ошибка ввода & ошибка ввода & F & -- & Пройден \\
\hline
3 & 0 & не простое & не простое & T & T & Пройден \\
\hline
\end{tabularx}
\end{table}

Все элементарные условия C1 и C2 проверены на оба логических значения. Условия C3--C6 аналогично покрыты тестами из таблицы покрытия решений.

\subsubsection{Покрытие решений и условий}

Критерий выполнен набором тестов, составленным для покрытия решений: каждое из условий C1--C6
принимает оба значения, все исходы ветвлений достигнуты хотя бы однократно.

\begin{table}[H]
\centering
\caption{Покрытие решений и условий}
\begin{tabularx}{\textwidth}{|c|c|X|X|c|c|c|c|c|c|c|}
\hline
\textbf{N} & \textbf{n} & \textbf{Ожид. рез.} & \textbf{Факт. рез.} & \textbf{C1} & \textbf{C2} & \textbf{C3} & \textbf{C4} & \textbf{C5} & \textbf{C6} & \textbf{Статус} \\
\hline
1 & 17 & простое & простое & T & F & F & F & T/F & F & Пройден \\
\hline
2 & \texttt{abc} & ошибка ввода & ошибка ввода & F & -- & -- & -- & -- & -- & Пройден \\
\hline
3 & 0 & не простое & не простое & T & T & -- & -- & -- & -- & Пройден \\
\hline
4 & 2 & простое & простое & T & F & T & -- & -- & -- & Пройден \\
\hline
5 & 8 & не простое & не простое & T & F & F & T & -- & -- & Пройден \\
\hline
6 & 15 & не простое & не простое & T & F & F & F & T & T & Пройден \\
\hline
\end{tabularx}
\end{table}

Все решения и условия покрыты.

\subsubsection{Комбинаторное покрытие условий}

Данный критерий требует проверки всех достижимых комбинаций значений элементарных условий.
Поскольку условия проверяются последовательно (при истинности предыдущего условие-прерыватель не 
вычисляет следующее), число достижимых комбинаций ограничено структурой программы.

\begin{enumerate}
    \item C1=F: программа немедленно запрашивает повторный ввод; C2--C6 недостижимы.
    \item C1=T, C2=T: программа выводит <<не простое>>; C3--C6 недостижимы.
    \item C1=T, C2=F, C3=T: программа выводит <<простое>>; C4--C6 недостижимы.
    \item C1=T, C2=F, C3=F, C4=T: программа выводит <<не простое>>; C5, C6 недостижимы.
    \item C1=T, C2=F, C3=F, C4=F, C5=T, C6=T: делитель найден $\to$ <<не простое>>.
    \item C1=T, C2=F, C3=F, C4=F, C5=T, C6=F: делитель не найден на данной итерации $\to$ следующая итерация.
    \item C1=T, C2=F, C3=F, C4=F, C5=F: цикл завершён без нахождения делителя $\to$ <<простое>>.
\end{enumerate}

\begin{table}[H]
\centering
\caption{Комбинаторное покрытие условий}
\begin{tabularx}{\textwidth}{|c|c|X|X|c|c|c|c|c|c|c|}
\hline
\textbf{N} & \textbf{n} & \textbf{Ожид. рез.} & \textbf{Факт. рез.} & \textbf{C1} & \textbf{C2} & \textbf{C3} & \textbf{C4} & \textbf{C5} & \textbf{C6} & \textbf{Статус} \\
\hline
1 & \texttt{abc} & ошибка ввода & ошибка ввода & F & -- & -- & -- & -- & -- & Пройден \\
\hline
2 & 0 & не простое & не простое & T & T & -- & -- & -- & -- & Пройден \\
\hline
3 & 2 & простое & простое & T & F & T & -- & -- & -- & Пройден \\
\hline
4 & 8 & не простое & не простое & T & F & F & T & -- & -- & Пройден \\
\hline
5 & 15 & не простое & не простое & T & F & F & F & T & T & Пройден \\
\hline
6 & 17 & простое & простое & T & F & F & F & T/F & F & Пройден \\
\hline
\end{tabularx}
\end{table}

Все достижимые комбинации условий покрыты. Тест №6 охватывает одновременно комбинации (5), (6) и (7), 
поскольку при $n=17$ цикл выполняет несколько итераций с C6=F, а затем C5 принимает значение F.

\subsubsection{Результаты тестирования программы №1 методом белого ящика}

В ходе тестирования методом белого ящика были разработаны тесты для пяти критериев покрытия кода.

Для покрытия операторов достаточно одного теста ($n = 17$): он проходит по пути
C1(T)$\to$C2(F)$\to$C3(F)$\to$C4(F)$\to$цикл$\to$вывод <<простое>> и охватывает все операторы программы.

Для покрытия решений потребовалось шесть тестов, которые в совокупности обеспечивают обе ветви
каждого из условий C1--C6. Аналогичный набор тестов удовлетворяет и критерию покрытия условий,
поскольку каждое из условий C1--C6 является элементарным.

Покрытие решений и условий достигается тем же набором из шести тестов.

Для комбинаторного покрытия выявлено семь достижимых комбинаций (все остальные комбинации
недостижимы из-за последовательного характера ветвлений). Шесть тестов покрывают все семь
достижимых комбинаций. Все тесты пройдены успешно.

\pagebreak

%=============================================================
\section{Тестирование программы №2}
%=============================================================

\subsection{Описание программы}

\noindent \textbf{Автор:} Попов Илья.

\noindent \textbf{Название:} алгоритм подсчета вхождений подстроки в массиве строк.

\noindent \textbf{Дано:}

\begin{enumerate}
	\item \texttt{arr} --- массив строк.
	\item \texttt{substr} --- подстрока для поиска в массиве.
\end{enumerate}

\noindent \textbf{Требуется:} найти количество вхождений \texttt{count} подстроки в массиве строк.
 
\noindent \textbf{Ограничения:}

\begin{enumerate}
	\item Длина $n$ подстроки \texttt{substr}: $n \in [0; 1000]$
	\item Длина $n$ строки \texttt{str} в массиве \texttt{arr}: $n \in [1; 1000]$
	\item Размер массива $n \in [1; 100]$
	\item \texttt{count} --- целое неотрицательное число.
\end{enumerate}

\noindent \textbf{Спецификация:}

\begin{table}[h!]
\centering
\begin{tabularx}{\textwidth}{|X|X|X|}
\hline
\textbf{Входные данные} & \textbf{Выходные данные} & \textbf{Реакция программы} \\
\hline
arr = \{"hello", "world", "hello"\} \newline substr = "he" &
count = 2 &
Подстрока встречается дважды. Возвращается значение. Программа завершает работу. \\
\hline
arr = \{"hello", "world", "hello"\} \newline substr = "" &
Error: "Подстрока не может быть пустой" &
Программа генерирует исключение. Завершение работы программы. \\
\hline
arr = \{"hello", "world", "hello"\} \newline substr = "eee" &
count = 0 &
Подстрока не встречается. Возвращается значение. Программа завершает работу. \\
\hline
arr = \{\} \newline substr = "he" &
Error: "Массив пустой" &
Программа генерирует исключение. Завершение работы программы. \\
\hline
длина arr > 100 \newline substr = "e" &
Error: "Массив слишком большой" &
Программа генерирует исключение. Завершение работы программы. \\
\hline
длина элемента в arr > 1000 \newline substr = "e" &
Error: "Слишком длинная строка в массиве" &
Программа генерирует исключение. Завершение работы программы. \\
\hline
\end{tabularx}
\end{table}

\begin{figure}[H]
	\centering
	\includegraphics[width=0.8\textwidth]{pics/scheme2.png}
	\captionof{figure}{Блок-схема программы №2}
\end{figure}

%-------------------------------------------------------------
\subsection{Тестирование методом <<чёрного>> ящика}
%-------------------------------------------------------------

\subsubsection{Разбиение на классы эквивалентности}

Исходя из ограничений программы составлено разбиение на классы эквивалентности
по двум входным параметрам: \texttt{arr} и \texttt{substr}.

{\centering\large{
\begin{longtable}{|p{4cm}|p{5.5cm}|p{5.5cm}|}
	\caption{Классы эквивалентности для программы №2}\\
	\hline
	Входное условие & Допустимые классы & Недопустимые классы \\ \hline
	Массив \texttt{arr} &
	Массив из 1 элемента (1) &
	Пустой массив (7) \\ \hline
	&
	Массив из 2--99 элементов (2) &
	Массив из более чем 100 элементов (8) \\ \hline
	&
	Массив ровно из 100 элементов (3) &
	Массив, содержащий строку длиной > 1000 (9) \\ \hline
	Подстрока \texttt{substr} &
	Непустая подстрока, встречающаяся в arr (4) &
	Пустая строка (10) \\ \hline
	&
	Непустая подстрока, не встречающаяся в arr (5) &
	Подстрока длиной > 1000 (11) \\ \hline
	&
	Подстрока длиной ровно 1000 (6) & \\ \hline
\end{longtable}}}

{\centering\large{
\begin{longtable}{|p{0.4cm}|p{4cm}|p{2.5cm}|p{3cm}|p{3cm}|p{2cm}|}
	\caption{Тесты по классам эквивалентности}\\
	\hline
	N & arr & substr & Ожидаемый результат & Фактический результат & Результат \\ \hline
	1 & \{"hello"\} & "he" & count = 1 & count = 1 & Пройден \\ \hline
	2 & \{"hello","world","hi"\} & "h" & count = 2 & count = 2 & Пройден \\ \hline
	3 & Массив из 100 элементов "a" & "a" & count = 100 & count = 100 & Пройден \\ \hline
	4 & \{"hello","world"\} & "world" & count = 1 & count = 1 & Пройден \\ \hline
	5 & \{"hello","world"\} & "xyz" & count = 0 & count = 0 & Пройден \\ \hline
	6 & \{"hello"\} & substr дл. 1000 & count = 0 & count = 0 & Пройден \\ \hline
	7 & \{\} & "he" & Ошибка: массив пустой & Ошибка: массив пустой & Пройден \\ \hline
	8 & Массив из 101 эл. & "a" & Ошибка: массив слишком большой & Ошибка: массив слишком большой & Пройден \\ \hline
	9 & \{"hello", str дл. 1001\} & "he" & Ошибка: строка слишком длинная & Ошибка: строка слишком длинная & Пройден \\ \hline
	10 & \{"hello"\} & "" & Ошибка: пустая подстрока & Ошибка: пустая подстрока & Пройден \\ \hline
	11 & \{"hello"\} & substr дл. 1001 & Ошибка: подстрока слишком длинная & Ошибка: подстрока слишком длинная & Пройден \\ \hline
\end{longtable}}}

Все 11 тестов прошли успешно.

\subsubsection{Анализ граничных условий}

Исходя из ограничений программы выделяются следующие граничные значения:

\begin{enumerate}
	\item Размер массива: 0 (пустой), 1 (минимальный допустимый), 100 (максимальный допустимый), 101 (первый недопустимый).
	\item Длина строки в массиве: 1 (минимальная), 1000 (максимальная), 1001 (первая недопустимая).
	\item Длина подстроки: 0 (пустая --- недопустима), 1 (минимальная допустимая), 1000 (максимальная), 1001 (первая недопустимая).
	\item Количество совпадений: 0 (нет вхождений), 1, все элементы.
\end{enumerate}

{\centering\large{
\begin{longtable}{|p{0.4cm}|p{3.5cm}|p{2.5cm}|p{3cm}|p{3cm}|p{2cm}|}
	\caption{Проверка граничных условий программы №2}\\
	\hline
	N & arr & substr & Ожидаемый результат & Фактический результат & Результат \\ \hline
	1 & \{\} & "a" & Ошибка: массив пустой & Ошибка: массив пустой & Пройден \\ \hline
	2 & \{"a"\} & "a" & count = 1 & count = 1 & Пройден \\ \hline
	3 & Массив из 100 эл. "a" & "a" & count = 100 & count = 100 & Пройден \\ \hline
	4 & Массив из 101 эл. & "a" & Ошибка: массив слишком большой & Ошибка: массив слишком большой & Пройден \\ \hline
	5 & \{"a"\} (дл. строки = 1) & "a" & count = 1 & count = 1 & Пройден \\ \hline
	6 & \{str дл. 1000\} & "a" & count = 1000 или 0 & соответствует & Пройден \\ \hline
	7 & \{str дл. 1001\} & "a" & Ошибка: строка слишком длинная & Ошибка: строка слишком длинная & Пройден \\ \hline
	8 & \{"hello"\} & "" & Ошибка: пустая подстрока & Ошибка: пустая подстрока & Пройден \\ \hline
	9 & \{"hello"\} & "h" (дл. 1) & count = 1 & count = 1 & Пройден \\ \hline
	10 & \{"hello"\} & substr дл. 1000 & count = 0 & count = 0 & Пройден \\ \hline
	11 & \{"hello"\} & substr дл. 1001 & Ошибка: подстрока слишком длинная & Ошибка: подстрока слишком длинная & Пройден \\ \hline
	12 & \{"a","b","c"\} & "x" & count = 0 & count = 0 & Пройден \\ \hline
	13 & \{"a","a","a"\} & "a" & count = 3 (все) & count = 3 & Пройден \\ \hline
\end{longtable}}}

Все 13 тестов прошли успешно.

\subsubsection{Причинно-следственная диаграмма}

Для построения причинно-следственной диаграммы введём следующие обозначения.

\textbf{Причины:}
\begin{enumerate}
    \item Массив \texttt{arr} пустой ($\texttt{len(arr)} = 0$).
    \item Подстрока \texttt{substr} пустая ($\texttt{len(substr)} = 0$).
    \item Массив слишком большой ($\texttt{len(arr)} > 100$).
    \item Подстрока слишком длинная ($\texttt{len(substr)} > 1000$).
    \item Строка в массиве слишком длинная ($\exists\, \texttt{str} \in \texttt{arr}: \texttt{len(str)} > 1000$).
    \item Подстрока встречается хотя бы в одном элементе массива ($\texttt{count} > 0$).
\end{enumerate}

\textbf{Следствия:}
\begin{enumerate}
    \setcounter{enumi}{20}
    \item Ошибка: <<Массив пустой>>.
    \item Ошибка: <<Подстрока не может быть пустой>>.
    \item Ошибка: <<Длина массива слишком большая>>.
    \item Ошибка: <<Длина подстроки слишком большая>>.
    \item Ошибка: <<Длина строки слишком большая>>.
    \item Вывод \texttt{count} $> 0$ (есть совпадения).
    \item Вывод \texttt{count} $= 0$ (нет совпадений).
\end{enumerate}

\textbf{Ограничения:}
\begin{itemize}
    \item E(1, 3): причины 1 и 3 взаимоисключающие (пустой массив не может одновременно быть слишком большим).
    \item При C1=T программа завершается до проверки C2--C6.
\end{itemize}

\begin{table}[H]
\centering
\caption{Таблица решений (причинно-следственная диаграмма)}
\renewcommand{\arraystretch}{1.4}
\begin{tabular}{|l|c|c|c|c|c|c|c|}
\hline
 & \textbf{1} & \textbf{2} & \textbf{3} & \textbf{4} & \textbf{5} & \textbf{6} & \textbf{7} \\
\hline
Причина 1 (arr пустой)         & 1 & 0 & 0 & 0 & 0 & 0 & 0 \\
\hline
Причина 2 (substr пустая)      & -- & 1 & 0 & 0 & 0 & 0 & 0 \\
\hline
Причина 3 (arr > 100)          & -- & -- & 1 & 0 & 0 & 0 & 0 \\
\hline
Причина 4 (substr > 1000)      & -- & -- & -- & 1 & 0 & 0 & 0 \\
\hline
Причина 5 (строка > 1000)      & -- & -- & -- & -- & 1 & 0 & 0 \\
\hline
Причина 6 (count $> 0$)        & -- & -- & -- & -- & -- & 1 & 0 \\
\hline
Следствие 21 (ошибка: arr пуст) & 1 & 0 & 0 & 0 & 0 & 0 & 0 \\
\hline
Следствие 22 (ошибка: substr пуста) & 0 & 1 & 0 & 0 & 0 & 0 & 0 \\
\hline
Следствие 23 (ошибка: arr > 100) & 0 & 0 & 1 & 0 & 0 & 0 & 0 \\
\hline
Следствие 24 (ошибка: substr > 1000) & 0 & 0 & 0 & 1 & 0 & 0 & 0 \\
\hline
Следствие 25 (ошибка: строка > 1000) & 0 & 0 & 0 & 0 & 1 & 0 & 0 \\
\hline
Следствие 26 (count $> 0$)     & 0 & 0 & 0 & 0 & 0 & 1 & 0 \\
\hline
Следствие 27 (count $= 0$)     & 0 & 0 & 0 & 0 & 0 & 0 & 1 \\
\hline
\end{tabular}
\end{table}

\begin{table}[H]
\centering
\caption{Тесты на основе таблицы решений}
\renewcommand{\arraystretch}{1.5}
\begin{tabularx}{\textwidth}{|c|X|c|X|X|c|}
\hline
\textbf{N} & \textbf{arr} & \textbf{substr} & \textbf{Ожидаемый рез.} & \textbf{Фактич. рез.} & \textbf{Результат} \\
\hline
1 & \{\} & "he" & Следствие 21: ошибка <<массив пустой>> & Следствие 21: ошибка <<массив пустой>> & Пройден \\
\hline
2 & \{"hello"\} & "" & Следствие 22: ошибка <<пустая подстрока>> & Следствие 22: ошибка & Пройден \\
\hline
3 & Массив из 101 эл. & "a" & Следствие 23: ошибка <<arr > 100>> & Следствие 23: ошибка & Пройден \\
\hline
4 & \{"hello"\} & substr дл. 1001 & Следствие 24: ошибка <<substr > 1000>> & Следствие 24: ошибка & Пройден \\
\hline
5 & \{"hello", str дл. 1001\} & "he" & Следствие 25: ошибка <<строка > 1000>> & Следствие 25: ошибка & Пройден \\
\hline
6 & \{"hello","world"\} & "hello" & Следствие 26: count = 1 & Следствие 26: count = 1 & Пройден \\
\hline
7 & \{"cat","dog"\} & "xyz" & Следствие 27: count = 0 & Следствие 27: count = 0 & Пройден \\
\hline
\end{tabularx}
\end{table}

Все 7 тестов прошли успешно.

\subsubsection{Результаты тестирования программы №2 методом чёрного ящика}

В результате тестирования методом чёрного ящика было составлено 31 тест
(11 тестов по классам эквивалентности + 13 граничных условий + 7 тестов по таблице решений).
Все тесты прошли успешно.

Тестирование по классам эквивалентности охватило все выделенные допустимые и недопустимые классы
для входных параметров \texttt{arr} и \texttt{substr}. Анализ граничных условий выявил корректность
обработки крайних значений по всем трём измерениям: размеру массива, длине строки и длине подстроки.
Тестирование на основе таблицы решений верифицировало каждое следствие в отдельности.

%-------------------------------------------------------------
\subsection{Тестирование методом белого ящика}
%-------------------------------------------------------------

Алгоритм составления тестов методом «белого» ящика предполагает обход всех
возможных путей в теле программы и проверку выполнения каждого оператора не
менее одного раза.

В программе можно выделить следующие ключевые условия и соответствующие им ветви:

\begin{enumerate}
    \item Условие C1: (\texttt{len(arr) == 0}):

        Ветвь T (True): Выход из программы с ошибкой «Массив пустой».

        Ветвь F (False): Продолжение выполнения.

    \item Условие C2: (\texttt{len(substr) == 0}):

        Ветвь T (True): Выход из программы с ошибкой «Подстрока не может быть пустой».

        Ветвь F (False): Продолжение выполнения.

    \item Условие C3: (\texttt{len(arr) > 100}):

        Ветвь T (True): Выход из программы с ошибкой «Длина массива слишком большая».

        Ветвь F (False): Продолжение выполнения.

    \item Условие C4: (\texttt{len(substr) > 1000}):

        Ветвь T (True): Выход из программы с ошибкой «Длина подстроки слишком большая».

        Ветвь F (False): Продолжение выполнения.

    \item Условие C5: (\texttt{len(str) > 1000}) (внутри цикла \texttt{foreach}):

        Ветвь T (True): Выход из программы с ошибкой «Длина строки слишком большая».

        Ветвь F (False): Продолжение проверки или выход из цикла.

    \item Условие C6: (\texttt{str.contains(substr)}):

        Ветвь T (True): Увеличение счётчика \texttt{count}.

        Ветвь F (False): Переход к следующему элементу.
\end{enumerate}

\noindent \textbf{Возможные пути выполнения:}

\begin{enumerate}
    \item Пустой массив: C1(T) $\to$ Завершение.
    \item Пустая подстрока: C1(F) $\to$ C2(T) $\to$ Завершение.
    \item Массив слишком большой: C1(F) $\to$ C2(F) $\to$ C3(T) $\to$ Завершение.
    \item Подстрока слишком длинная: C1(F) $\to$ C2(F) $\to$ C3(F) $\to$ C4(T) $\to$ Завершение.
    \item Строка в массиве слишком длинная: C1(F) $\to$ \ldots $\to$ C4(F) $\to$ цикл $\to$ C5(T) $\to$ Завершение.
    \item Успешный поиск (есть совпадения): C1(F) $\to$ \ldots $\to$ C4(F) $\to$ цикл $\to$ C5(F) $\to$ C6(T) $\to$ \ldots $\to$ Вывод \texttt{count}.
    \item Поиск без совпадений: C1(F) $\to$ \ldots $\to$ C4(F) $\to$ цикл $\to$ C5(F) $\to$ C6(F) $\to$ \ldots $\to$ Вывод \texttt{count = 0}.
    \item Смешанный поиск: C1(F) $\to$ \ldots $\to$ C4(F) $\to$ цикл с многократным C5(F), C6(T/F) $\to$ Вывод \texttt{count}.
\end{enumerate}

\subsubsection{Покрытие операторов}

Для покрытия всех операторов достаточно одного теста, проходящего по пути с успешным поиском.

\begin{table}[h!]
\centering
\caption{Покрытие операторов}
\begin{tabularx}{\textwidth}{|c|c|X|X|c|c|}
\hline
\textbf{N} & \textbf{arr} & \textbf{substr} & \textbf{Ожид. результат} & \textbf{Факт. результат} & \textbf{Статус} \\
\hline
1 & \{"hello","world","hello"\} & "he" & 2 & 2 & Пройден \\
\hline
\end{tabularx}
\end{table}

Тест №1 выполнил все операторы программы.

\subsubsection{Покрытие решений}

\begin{table}[h!]
\centering
\caption{Покрытие решений}
\begin{tabularx}{\textwidth}{|c|X|X|X|X|c|c|c|c|c|c|c|}
\hline
\textbf{N} & \textbf{arr} & \textbf{substr} & \textbf{Ожид. рез.} & \textbf{Факт. рез.} & \textbf{C1} & \textbf{C2} & \textbf{C3} & \textbf{C4} & \textbf{C5} & \textbf{C6} & \textbf{Статус} \\
\hline
1 & \{"hello","world","hello"\} & "he" & 2 & 2 & F & F & F & F & F & T & Пройден \\
\hline
2 & \{"cat","dog"\} & "xyz" & 0 & 0 & F & F & F & F & F & F & Пройден \\
\hline
3 & \{\} & "he" & Ошибка: массив пустой & Ошибка: массив пустой & T & -- & -- & -- & -- & -- & Пройден \\
\hline
4 & \{"hello"\} & "" & Ошибка: пустая подстрока & Ошибка: пустая подстрока & F & T & -- & -- & -- & -- & Пройден \\
\hline
5 & arr из 101 эл. & "a" & Ошибка: массив слишком большой & Ошибка: массив слишком большой & F & F & T & -- & -- & -- & Пройден \\
\hline
6 & \{"hello"\} & substr дл. 1001 & Ошибка: подстрока слишком длинная & Ошибка: подстрока слишком длинная & F & F & F & T & -- & -- & Пройден \\
\hline
7 & \{"hello", str дл. 1001\} & "he" & Ошибка: строка слишком длинная & Ошибка: строка слишком длинная & F & F & F & F & T & -- & Пройден \\
\hline
\end{tabularx}
\end{table}

Все условия были протестированы на T и F.

\subsubsection{Покрытие условий}

\begin{table}[h!]
\centering
\caption{Покрытие условий}
\begin{tabularx}{\textwidth}{|c|X|X|X|X|c|c|c|}
\hline
\textbf{N} & \textbf{arr} & \textbf{substr} & \textbf{Ожид. рез.} & \textbf{Факт. рез.} & \textbf{C1} & \textbf{C2} & \textbf{Статус} \\
\hline
1 & \{"hello","world"\} & "he" & 1 & 1 & F & F & Пройден \\
\hline
2 & \{\} & "he" & Ошибка: массив пустой & Ошибка: массив пустой & T & -- & Пройден \\
\hline
3 & \{"hello"\} & "" & Ошибка: пустая подстрока & Ошибка: пустая подстрока & F & T & Пройден \\
\hline
\end{tabularx}
\end{table}

Все элементарные условия были проверены на T и F.

\subsubsection{Покрытие решений и условий}

Тесты, написанные для покрытия решений, обеспечивают и покрытие решений и условий.

\begin{table}[h!]
\centering
\caption{Покрытие решений и условий}
\begin{tabularx}{\textwidth}{|c|X|X|X|X|c|c|c|c|c|c|c|}
\hline
\textbf{N} & \textbf{arr} & \textbf{substr} & \textbf{Ож. рез.} & \textbf{Факт. рез.} & \textbf{C1} & \textbf{C2} & \textbf{C3} & \textbf{C4} & \textbf{C5} & \textbf{C6} & \textbf{Статус} \\
\hline
1 & \{"hello","world","hello"\} & "he" & 2 & 2 & F & F & F & F & F & T & Пройден \\
\hline
2 & \{"cat","dog"\} & "xyz" & 0 & 0 & F & F & F & F & F & F & Пройден \\
\hline
3 & \{\} & "he" & Ошибка: массив пустой & Ошибка: массив пустой & T & -- & -- & -- & -- & -- & Пройден \\
\hline
4 & \{"hello"\} & "" & Ошибка: пустая подстрока & Ошибка: пустая подстрока & F & T & -- & -- & -- & -- & Пройден \\
\hline
5 & arr из 101 эл. & "a" & Ошибка: массив слишком большой & Ошибка: массив слишком большой & F & F & T & -- & -- & -- & Пройден \\
\hline
6 & \{"hello"\} & substr дл. 1001 & Ошибка: подстрока слишком длинная & Ошибка: подстрока слишком длинная & F & F & F & T & -- & -- & Пройден \\
\hline
7 & \{"hello", str дл. 1001\} & "he" & Ошибка: строка слишком длинная & Ошибка: строка слишком длинная & F & F & F & F & T & -- & Пройден \\
\hline
\end{tabularx}
\end{table}

Все решения и условия покрыты.

\subsubsection{Комбинаторное покрытие условий}

Для условий C1 и C2 возможны комбинации (F,F), (T,--), (F,T). Комбинация (T,T) недостижима,
так как при C1=T программа завершается до проверки C2.

\begin{table}[H]
\centering
\caption{Комбинаторное покрытие условий}
\begin{tabularx}{\textwidth}{|c|X|X|X|X|c|c|c|}
\hline
\textbf{N} & \textbf{arr} & \textbf{substr} & \textbf{Ожид. рез.} & \textbf{Факт. рез.} & \textbf{C1} & \textbf{C2} & \textbf{Статус} \\
\hline
1 & \{"hello","world","hello"\} & "he" & 2 & 2 & F & F & Пройден \\
\hline
2 & \{\} & "he" & Ошибка: массив пустой & Ошибка: массив пустой & T & -- & Пройден \\
\hline
3 & \{"hello"\} & "" & Ошибка: пустая подстрока & Ошибка: пустая подстрока & F & T & Пройден \\
\hline
4 & arr из 101 эл. & "a" & Ошибка: массив слишком большой & Ошибка: массив слишком большой & F & F (C3=T) & Пройден \\
\hline
5 & \{"hello"\} & substr дл. 1001 & Ошибка: подстрока слишком длинная & Ошибка: подстрока слишком длинная & F & F (C4=T) & Пройден \\
\hline
6 & \{"hello", str дл. 1001\} & "he" & Ошибка: строка слишком длинная & Ошибка: строка слишком длинная & F & F (C5=T) & Пройден \\
\hline
7 & \{"cat","dog"\} & "xyz" & 0 & 0 & F & F (C6=F) & Пройден \\
\hline
\end{tabularx}
\end{table}

Все достижимые комбинации условий покрыты. Комбинация C1=T и C2=T недостижима.

\subsubsection{Результаты тестирования программы №2 методом белого ящика}

В ходе тестирования алгоритма подсчёта вхождений подстроки в массиве строк методом белого ящика
были разработаны тесты для пяти критериев покрытия кода. Все тесты пройдены успешно, ошибок в
логике программы обнаружено не было. Подробные результаты приведены в таблицах выше.

\pagebreak

\section{Заключение}

В ходе выполнения данной лабораторной работы были изучены и применены методы тестирования
«белым ящиком» и «чёрным ящиком» для обеих программ.

\textbf{Программа №1 (проверка простоты числа).}
Тестирование методом «чёрного ящика» включало 26 тестов: 12 по классам эквивалентности,
8 по граничным значениям и 6 по таблице решений. Все тесты пройдены успешно.
Тестирование методом «белого ящика» включало 6 тестов, охватывающих пять критериев покрытия
(операторов, решений, условий, решений и условий, комбинаторное). Все 7 достижимых комбинаций
условий C1--C6 покрыты. Дефектов обнаружено не было.

\textbf{Программа №2 (подсчёт вхождений подстроки).}
Тестирование методом «чёрного ящика» включало 31 тест: 11 по классам эквивалентности,
13 по граничным значениям и 7 по таблице решений. Все тесты пройдены успешно.
Тестирование методом «белого ящика» включало 7 тестов, охватывающих пять критериев покрытия
кода. Комбинация C1=T, C2=T признана недостижимой вследствие раннего завершения программы
при C1=T. Все тесты пройдены успешно.

По итогам выполненной работы можно сделать вывод, что оба метода тестирования дополняют
друг друга и в совокупности обеспечивают высокое качество проверки программного обеспечения.
Метод «белого ящика» позволяет убедиться в корректности внутренней логики и полноте охвата
всех ветвей кода. Метод «чёрного ящика» проверяет соответствие поведения программы требованиям
спецификации без учёта внутренней реализации. Совместное применение обоих методов снижает
вероятность пропуска дефектов и обеспечивает комплексную верификацию программного продукта.

\pagebreak

\begin{thebibliography}{3}
	\bibitem{[1]} Майерс Г. <<Искусство тестирования программ>>, СПб, Диалектика, 2012г.
\end{thebibliography}

\end{document}